% !TEX root = ../SANER2027-main-labeling.tex
\section{Experimental Setup}
\label{sec:setup}

\myparagraph{Dataset and Settings}
\label{sec:setup-dataset}
We use the dataset introduced in Heo et al.~\cite{heo2025study} for our experiments. The dataset has a total of 6,600 issue reports gathered from eleven active open-source projects:
\texttt{ansible}, \texttt{bitcoin},
\texttt{dart-lang}, \texttt{roslyn},
\texttt{react}, \texttt{flutter},
\texttt{kubernetes}, \texttt{TypeScript},
\texttt{vscode}, \texttt{opencv}, and
\texttt{tensorflow}. Each project has 600 issue reports
with an equal distribution across the labels
(\emph{bug}, \emph{feature}, and \emph{question}). To support their LLM fine-tuning method, Heo et al. divided the dataset into 50/50 train and test
splits, stratified by project and label. This results in 300 train
and 300 test issue reports per project. We also adopt Heo et al.'s project-specific (PS) and project-agnostic (PA) configurations to examine how data scope affects IRC performance. In the PS setting, each project's train and test data are used in isolation, resulting in eleven different train-test pairs. In the PA setting, the train splits of all projects are merged into a single training set of 3,300 issue reports. In both configurations, \rag and \knn only use the training splits as their retrieval index to ensure fair comparison with the fine-tuning baseline.

%%%%%%%%%%%%
\myparagraph{Models}
\label{sec:setup-models}
We use the Qwen2.5-Instruct model family~\cite{yang2024qwen25} in our evaluations with four different parameter sizes: 3B, 7B, 14B, and 32B. Qwen2.5 is an open-source LLM with strong performance on general-purpose benchmarks. Using the same model with different parameter sizes enables analysis of the effect of model scale in LLM-based IRC approaches.


%%%%%%%%%%%%
\myparagraph{Evaluation Metrics}
\label{sec:setup-metrics}
We report the macro-averaged $F_1$ over the three labels as the headline metric, since it weighs the three labels equally. A method's \emph{best} $k$ (its best configuration) is the one with the highest macro $F_1$ on the test set. We report scores in percent and differences between methods in percentage points (points). We also report macro-averaged and per-class precision and recall. All metrics are calculated over the total 3{,}300-issue test set. We also report the rate of invalid outputs (\textit{invalid rate}), which count as incorrect predictions. For significance, we report bootstrap 95\% confidence intervals (CIs) on the macro $F_1$ differences between methods.\footnote{We resample the 3{,}300 test issues 1{,}000 times and evaluate both methods on the same resample (paired bootstrap).}

%%%%%%%%%%%%
\myparagraph{Hardware Setup}
\label{sec:setup-hardware}
Fine-tuning of Qwen-14B and Qwen-32B and the Qwen-32B runs at $k{\ge}12$ used one NVIDIA L40 or L40S GPU (48~GB VRAM). All other experiments used one NVIDIA RTX 4090 GPU (24~GB VRAM).
